\section{Summary and outlook}
\label{sec:conclusion}
NNPDF3.1  is the new main PDF release from the
NNPDF family.
%
It represents a significant improvement over
NNPDF3.0, by including constraints from many new observables,
some of which are included for the first time in a global
PDF determination,  thanks to the recent availability of the corresponding
NNLO QCD corrections. Notable examples are $t\bar{t}$ differential distributions
and the $Z$ boson $p_T$ spectrum.
%
From the theory point of view,
the main improvement is to place the charm
PDF on an equal footing as the light quark
PDFs.
%
Independently parametrizing  the charm PDF  resolves a tension which would otherwise be present
between ATLAS gauge boson production and HERA inclusive structure
function data, leads to improved agreement with the 
LHC data, and turns the strong dependence of perturbatively
generated charm on the value of the pole charm mass into a PDF
uncertainty, as most of the mass dependence is reabsorbed into the initial
PDF shape. 

The NNPDF3.1 set is also the first set for which PDFs are delivered in a
variety of formats: first of all, they are released both in Hessian
and Monte Carlo form, and furthermore, the default sets are optimized
and compressed so that a smaller number of  
Monte Carlo replicas or Hessian error sets reproduces the statistical
features of much larger underlying replica sets.
We now   discuss how both Hessian and Monte
Carlo reduced sets have
been produced out of a large set of Monte Carlo
replicas;  we then summarize all PDF sets that have been made public
through the LHAPDF interface; and  finally  we
present a brief outlook on future developments.

\subsection{Validation of the NNPDF3.1 reduced sets}
\label{sec:compression}

  Default NNPDF3.1 NLO and NNLO PDFs for the central
  $\alpha_s(m_Z)=0.118$ value, as well as the modified version with
  perturbative charm discussed in Sect.~\ref{sec:results-mc}, have been produced as  $N_{\rm
    rep}=1000$ replica sets. These large replica samples have been
subsequently processed using two reduction
  strategies: the Compressed Monte Carlo (CMC) algorithm~\cite{Carrazza:2015hva},
  to obtain
  a Monte Carlo representation based on a smaller number of replicas,
  and the  MC2H algorithm~\cite{Carrazza:2015aoa},
  to achieve an optimal Hessian representation
  of the underlying PDF probability distribution with a fixed number
  of error sets. Specifically,  we have thus constructed
  CMC-PDF sets with $N_{\rm rep}=100$ replicas and
  MC2H sets with $N_{\rm eig}=100$ (symmetric) eigenvectors.

  In Fig.~\ref{fig:cmc-mc2h-validation} we show the
  comparison between the PDFs from the input set
       of $N_{\rm rep}=1000$ replicas of NNPDF3.1 NNLO with the corresponding
       reduced sets of the CMC-PDFs with $N_{\rm rep}=100$ replicas
       and the MC2H hessian PDFs with $N_{\rm eig}=100$
       eigenvalues.
       %
       The agreement between the input $N_{\rm rep}=1000$ replica MC
       PDFs and the two reduced sets is very good in all cases.
       %
       By construction, the agreement in central values
       and PDF variances is slightly better for the MC2H sets,
       since the CMC-PDF sets aim to reproduce also higher
       moments in the probability distribution and thus possibly
       non-gaussian features, while Hessian sets are Gaussian by
       construction. 
       %
       Following the analysis of~\cite{Carrazza:2015aoa,Carrazza:2015hva}
       we have verified that also the correlations between PDFs
       are reproduced to a high degree of accuracy.
       
%%%%%%%%%%%%%%%%%%%%%%%%%%%%%%%%%%%%%%%%%%%%%%%%%%%%%%%%%%%%%%%%%%%%%
\begin{figure}[t]
  \begin{center}
     \includegraphics[scale=0.32]{images/plots/xu-cmc-mc2h-nnlo-fc.pdf}
     \includegraphics[scale=0.32]{images/plots/xd-cmc-mc2h-nnlo-fc.pdf}
     \includegraphics[scale=0.32]{images/plots/xubar-cmc-mc2h-nnlo-fc.pdf}
     \includegraphics[scale=0.32]{images/plots/xdbar-cmc-mc2h-nnlo-fc.pdf}
     \includegraphics[scale=0.32]{images/plots/xsp-cmc-mc2h-nnlo-fc.pdf}
     \includegraphics[scale=0.32]{images/plots/xc-cmc-mc2h-nnlo-fc.pdf}
     \includegraphics[scale=0.32]{images/plots/xsinglet-cmc-mc2h-nnlo-fc.pdf}
     \includegraphics[scale=0.32]{images/plots/xg-cmc-mc2h-nnlo-fc.pdf}
     \caption{\small Comparison between the PDFs from the input set
       of $N_{\rm rep}=1000$ replicas of NNPDF3.1 NNLO, the 
       reduced Monte Carlo CMC-PDFs with $N_{\rm rep}=100$ replicas,
       and the MC2H hessian PDFs with $N_{\rm eig}=100$ symmetric
       eigenvalues.
    \label{fig:cmc-mc2h-validation}}
\end{center}
\end{figure}
%%%%%%%%%%%%%%%%%%%%%%%%%%%%%%%%%%%%%%%%%%%%%%%%%%%%%%%%%%%%%%%%%%%%%%

In order to validate the efficiency of the CMC-PDF algorithm reduction
from the starting 
$N_{\rm rep}=1000$ replicas down to the compressed $N_{\rm rep}=100$ replicas,
in Fig.~\ref{fig:cmc-mc2h-validation-nrep} we show, following the procedure
described in~\cite{Carrazza:2015hva}, the
summary of statistical estimators
       that compare specific properties of the probability distributions defined
       by the input $N_{\rm rep}=1000$ replicas of NNPDF3.1 NNLO and the corresponding
       compressed sets as a function of $\widetilde{N}_{\rm rep}$, the number
       of replicas in the reduced set starting from $\widetilde{N}_{\rm rep}=100$.
       %
       We compare the results of the compression algorithm with those
       of random selection of  $\widetilde{N}_{\rm rep}$ replicas out of the original
       1000 ones: the error function ERF corresponding to central
       values, standard 
       deviations, kurtosis and skewness, correlations and the
       Kolmogorov distance are all shown. 
       %
       These results indicate that a CMC-PDF  100 replica set
       reproduces roughly the information contained in a random 
       $\widetilde{N}_{\rm rep}=400$  PDF set.


%%%%%%%%%%%%%%%%%%%%%%%%%%%%%%%%%%%%%%%%%%%%%%%%%%%%%%%%%%%%%%%%%%%%%
\begin{figure}[t]
  \begin{center}
    \includegraphics[scale=0.35]{images/plots/cmc_cv.pdf}\includegraphics[scale=0.35]{images/plots/cmc_std.pdf}
    \includegraphics[scale=0.35]{images/plots/cmc_skewness.pdf}\includegraphics[scale=0.35]{images/plots/cmc_kurtosis.pdf}
    \includegraphics[scale=0.35]{images/plots/cmc_kolmogorov.pdf}\includegraphics[scale=0.35]{images/plots/cmc_correlation.pdf}

     \caption{\small Comparison of estimators
       of the probability distributions computed using the NNPDF3.1
       NNLO input $N_{\rm
         rep}=1000$ replica set, and compressed sets of
       $\widetilde{N}_{\rm rep}$ replicas, plotted as a function of  
 $\widetilde{N}_{\rm rep}$.
       The error function (ERF) corresponding to central values, standard
       deviations, kurtosis, skewness, correlations and 
       Kolmogorov distance are shown.
    \label{fig:cmc-mc2h-validation-nrep}
  }
\end{center}
\end{figure}
%%%%%%%%%%%%%%%%%%%%%%%%%%%%%%%%%%%%%%%%%%%%%%%%%%%%%%%%%%%%%%%%%%%%%%

\clearpage



\subsection{Delivery}
\label{sec:delivery}

We now  provide a full list of the NNPDF3.1 PDF sets that are being 
made publicly available via the {\sc\small LHAPDF6}
interface~\cite{Buckley:2014ana}, 
\begin{center}
{\bf \url{http://lhapdf.hepforge.org/}~} .
\end{center}
As repeatedly mentioned in the paper, a very wide  set of results
concerning these PDF sets is available from the repository
\begin{center}
{\bf \url{http://nnpdf.hepforge.org/html/nnpdf31/catalog}~} .
\end{center}
All sets are made available as  $N_{\rm rep}=100$ Monte Carlo sets. 
For the baseline
sets, these are constructed out of larger  $N_{\rm rep}=1000$ replica
sets, which are also being made available. The baseline sets are also
provided as Hessian sets with 100 error sets.

The full list is the following:

\begin{itemize}

\item {\it Baseline NLO and NNLO NNPDF3.1 sets}

Baseline NLO and NNLO NNPDF3.1 sets are based on the global
dataset, with $\alpha_s(m_Z)=0.118$ and a variable-flavor number with
up to five active flavors. These sets contain $N_{\rm rep}=1000$ PDF replicas. 
\begin{flushleft}
\tt NNPDF31\_nlo\_as\_0118\_1000 \\
\tt NNPDF31\_nnlo\_as\_0118\_1000.\\
\end{flushleft}
A modified version in which charm is perturbatively generated (ad in
previous NNPDF sets) is also being made available, also with $N_{\rm rep}=1000$ PDF replicas:
\begin{flushleft}
\tt NNPDF31\_nlo\_pch\_as\_0118\_1000 \\
\tt NNPDF31\_nnlo\_pch\_as\_0118\_1000.\\
\end{flushleft}
Out of these , optimized  Monte Carlo $N_{\rm rep}=100$ 
\begin{flushleft}
\tt NNPDF31\_nlo\_as\_0118 \\
\tt NNPDF31\_nnlo\_as\_0118 \\
\tt NNPDF31\_nlo\_pch\_as\_0118 \\
\tt NNPDF31\_nnlo\_pch\_as\_0118,
\end{flushleft}
and  Hessian sets with $N_{\rm eig}=100$ eigenvectors 
\begin{flushleft}
\tt NNPDF31\_nlo\_as\_0118\_hessian \\
\tt NNPDF31\_nnlo\_as\_0118\_hessian \\
\tt NNPDF31\_nlo\_pch\_as\_0118\_hessian \\
\tt NNPDF31\_nnlo\_pch\_as\_0118\_hessian
\end{flushleft}
have been constructed as discussed in Sect.~\ref{sec:compression}
above. 

Out of these, smaller sets of eigenvectors optimized for the
computation of specific observables may be constructed using the {\tt
  SM-PDF} tool~\cite{Carrazza:2016htc}.
Specifically, sets optimized for a wide list of predefined observables
can be generated and downloaded using the web
interface~\cite{Carrazza:2016wte}  at
\begin{center}
{\bf \url{https://smpdf.mi.infn.it}~} .
\end{center}





\item {\it Flavor number variation}

We have produced sets, both at NLO and NNLO 
in which the maximum number of flavors differs
from the default 5, and it is either extended up to six, or frozen at
four:
\begin{flushleft}
  \tt NNPDF31\_nlo\_as\_0118\_nf\_4 \\
  \tt NNPDF31\_nlo\_as\_0118\_nf\_6 \\
  \tt NNPDF31\_nnlo\_as\_0118\_nf\_4 \\
  \tt NNPDF31\_nnlo\_as\_0118\_nf\_6. \\
\end{flushleft}
The variant with perturbatively generated charm is also made
available, in this case also in a $n_f=3$ fixed-flavor number scheme:
\begin{flushleft}
  \tt NNPDF31\_nlo\_pch\_as\_0118\_nf\_3 \\
  \tt NNPDF31\_nlo\_pch\_as\_0118\_nf\_4 \\
  \tt NNPDF31\_nlo\_pch\_as\_0118\_nf\_6 \\
  \tt NNPDF31\_nnlo\_pch\_as\_0118\_nf\_3 \\
  \tt NNPDF31\_nnlo\_pch\_as\_0118\_nf\_4 \\
  \tt NNPDF31\_nnlo\_pch\_as\_0118\_nf\_6. \\
\end{flushleft}




\item {\it $\alpha_s$ variation.}

We have produced NNLO sets with the following values of $\alpha_s(m_Z)$
0.108, 1.110, 0.112, 0.114, 0.116, 0.117, 0.118, 0.119, 0.120, 0.122,
0.124:
\begin{flushleft}
  \tt NNPDF31\_nnlo\_as\_0108 \\
  \tt NNPDF31\_nnlo\_as\_0110 \\
  \tt NNPDF31\_nnlo\_as\_0112 \\
  \tt NNPDF31\_nnlo\_as\_0114 \\
  \tt NNPDF31\_nnlo\_as\_0116 \\
  \tt NNPDF31\_nnlo\_as\_0117 \\
  \tt NNPDF31\_nnlo\_as\_0118 \\
  \tt NNPDF31\_nnlo\_as\_0119 \\
  \tt NNPDF31\_nnlo\_as\_0120 \\
  \tt NNPDF31\_nnlo\_as\_0122 \\
  \tt NNPDF31\_nnlo\_as\_0124. \\
\end{flushleft}

For the values $\alpha_s(m_Z)=0.116$ and $\alpha_s(m_Z)=0.120$ we have
also produced  NLO sets, 
\begin{flushleft}
  \tt NNPDF31\_nlo\_as\_0116 \\
  \tt NNPDF31\_nlo\_as\_0120, \\
%  \tt NNPDF31\_nnlo\_as\_0116 \\
%  \tt NNPDF31\_nnlo\_as\_0120,\\
\end{flushleft}
and the  variant with perturbative charm both at NLO and NNLO
\begin{flushleft}
  \tt NNPDF31\_nlo\_pch\_as\_0116 \\
  \tt NNPDF31\_nlo\_pch\_as\_0120 \\
\tt NNPDF31\_nnlo\_pch\_as\_0116\\
\tt NNPDF31\_nnlo\_pch\_as\_0120.\\
\end{flushleft}


In order to facilitate the computation of the combined PDF+$\alpha_s$
uncertainties we have also provided bundled  PDF+$\alpha_s$ variation
sets for $\alpha_s(m_Z)=0.118\pm0.002$. These are
provided both as Monte Carlo sets, and as Hessian sets. 
\begin{flushleft}
  \tt
NNPDF31\_nlo\_pdfas	\\
NNPDF31\_nnlo\_pdfas	\\
NNPDF31\_nlo\_pch\_pdfas \\	
NNPDF31\_nnlo\_pch\_pdfas  \\	
NNPDF31\_nlo\_hessian\_pdfas \\
NNPDF31\_nnlo\_hessian\_pdfas \\	
NNPDF31\_nlo\_pch\_hessian\_pdfas \\	
NNPDF31\_nnlo\_pch\_hessian\_pdfas. \\
\end{flushleft}
They are constructed as follows:
\begin{enumerate}
\item The central value (PDF member 0) is the central
  value of the corresponding $\alpha_s(m_Z)=0.118$ set.
\item The PDF members 1 to 100 correspond to the
  $N_{\rep}=100$ ($N_{\rm eig}=100$) Monte Carlo replicas
  (Hessian eigenvectors) from the $\alpha_s(m_Z)=0.118$ set.
\item The PDF members 101 and 102 are the central values
  of the sets with $\alpha_s(m_Z)=0.116$ and $\alpha_s(m_Z)=0.120$
  respectively.
\end{enumerate}
Note that, therefore, in the Hessian case member 0 is the central set,
and all remaining  bundled members $1-102$ are error sets, while in the Monte
Carlo case,  members $1-100$ are Monte Carlo replicas, while members
0,~$101$ and~$102$ are central sets (replica averages). The way they
should be used to compute combined PDF+$\alpha_s$ uncertainties is
discussed e.g. in Ref.~\cite{Butterworth:2015oua}.


\item {\it Charm mass variation}


We  provide sets with different values of the
  charm mass $m_c^{\rm pole}$.
  %
  They are available only at NNLO
with $m_c^{\rm pole}=1.38$ GeV and $m_c^{\rm pole}=1.64$ GeV. 
  \begin{flushleft}
    {\tt NNPDF31\_nnlo\_as\_0118\_mc\_138}\\
    {\tt NNPDF31\_nnlo\_as\_0118\_mc\_164}.\\
\end{flushleft}
For
comparison, the corresponding modified version with perturbative
charm are also made available:
  \begin{flushleft}
     {\tt NNPDF31\_nnlo\_pch\_as\_0118\_mc\_138}\\
    {\tt NNPDF31\_nnlo\_pch\_as\_0118\_mc\_164}.\\
\end{flushleft}


\item {\it Forced positivity sets}

We provide sets in which PDFs are non-negative:
  \begin{flushleft}
    {\tt NNPDF31\_nlo\_as\_0118\_mc}\\
     {\tt NNPDF31\_nlo\_pch\_as\_0118\_mc}\\
    {\tt NNPDF31\_nnlo\_as\_0118\_mc}\\
     {\tt NNPDF31\_nnlo\_pch\_as\_0118\_mc}\\
\end{flushleft}
These have been constructed simply setting to zero PDFs whenever they
become negative. They are thus an approximation, provided for
convenience for use in conjunction with codes which fail when PDFs are
negative.

\item {\it LO sets} 

Leading-order PDF sets are made available 
 $\alpha_s = 0.118$ and  $\alpha_s = 0.130$: 
  \begin{flushleft}
    {\tt NNPDF31\_lo\_as\_0118}\\
    {\tt NNPDF31\_lo\_as\_0130}\\
\end{flushleft}
The corresponding variant with perturbative charm is also provided:
  \begin{flushleft}
    {\tt NNPDF31\_lo\_pch\_as\_0118}\\
    {\tt NNPDF31\_lo\_pch\_as\_0130}.
\end{flushleft}


\item {\it Reduced datasets}

PDFs determined from subsets of the full NNPDF3.1, discussed in
  Sect.~\ref{sec:impactzpt}-\ref{sec:collideronly} are also made
  available, specifically
  \begin{flushleft}
  \begin{tabular}{ll}
  {\tt NNPDF31\_nnlo\_as\_0118\_noZpt} & no $Z$ $p_T$ data, see Sect.~\ref{sec:impactzpt}; \\
  {\tt NNPDF31\_nnlo\_as\_0118\_notop} & no $t\bar{t}$ production data, see Sect.~\ref{sec:impacttop}; \\
  {\tt NNPDF31\_nnlo\_as\_0118\_nojets} & no jet data,  see Sect.~\ref{sec:jetdata}; \\
  {\tt NNPDF31\_nnlo\_as\_0118\_wEMC} & EMC charm data added, see Sect.~\ref{sec:emc};\\
  {\tt NNPDF31\_nnlo\_as\_0118\_noLHC} & no LHC data,  see Sect.~\ref{sec:nolhc};  \\
  {\tt NNPDF31\_nnlo\_as\_0118\_proton}  &  proton-target data only, see Sect.~\ref{sec:nonucl};\\
  {\tt NNPDF31\_nnlo\_as\_0118\_collider} & collider data only, see Sect.~\ref{sec:collideronly}.\\
  \end{tabular}
  \end{flushleft}
\end{itemize}

 In addition to these, any  other PDF sets  discussed in this paper  is also available upon request.

\subsection{Outlook}
\label{sec:outlook}

The NNPDF3.1 PDF determination presented here is an update of the
previous NNPDF3.0, yet it contains substantial innovations both in
terms of methodology and dataset and it leads to  substantially more
precise and accurate PDF sets. Thanks to this, several spin-offs can
be pursued, either with the goal of updating existing results
based on previous NNPDF releases, or in some cases because the greater
accuracy enables projects which  previously were either impossible or
uninteresting.

These spin-off projects include the following:
\begin{itemize}
\item A precision
determination of the strong coupling constant $\alpha_s(m_Z)$.
%
We expect a significant increase in both precision and accuracy compared to previous
NNPDF determinations~\cite{Lionetti:2011pw,Ball:2011us}. Specifically,
thanks to the inclusion of many collider observables
at NNLO with direct dependence on both the gluon and on
$\alpha_s(m_Z)$
 it might turn out to be advantageous to drop 
altogether data taken on nuclear targets, and at low scales, {\it i.e.} base
the $\alpha_s$ determination on the  collider-only dataset of
Sect.~\ref{sec:collideronly}, or possibly an even more conservative dataset.
\item A determination of the charm mass  $m_c$. This would be for the
  first time based on a PDF determination in which the charm PDF is
  independently parametrized, thereby avoiding bias related to the identification of $m_c$
  as a parameter which determines the size of the charm PDF.
  %
\item A NNPDF3.1QED PDF set including the photon PDF $\gamma(x,Q^2)$, thereby updating
  the previous NNPDF2.3QED~\cite{Ball:2013hta} and
  NNPDF3.0QED~\cite{Bertone:2016ume} sets. This update should
  include recent theoretical progress: specifically the direct ``LuxQED''
  constraints which determine  the photon PDF from nucleon
  structure functions~\cite{Manohar:2016nzj};
and  NLO QED corrections to PDF evolution and DIS coefficient
functions, now included in {\tt APFEL}~\cite{Bertone:2013vaa}. 
\item PDF sets including small-$x$ resummation, based on the formalism
developed in~\cite{Altarelli:2005ni,Ball:2007ra,Altarelli:2008aj}, which
has been implemented in the {\tt HELL} code~\cite{Bonvini:2016wki}
and interfaced to {\tt APFEL}.
%
These  would provide an  answer the long-standing
issue of whether or not small-$x$ inclusive HERA data are adequately described
by fixed-order perturbative
evolution~\cite{Caola:2010cy,Caola:2009iy}, and may ultimately give
better control  
of theoretical uncertainties at small $x$.
\end{itemize}

On a longer timescale, further substantial improvements in dataset and
methodology are expected.
On the one hand, so far we have essentially  
restricted ourselves to LHC 8 TeV data.
A future release will include a significant
number of 13 TeV measurements, of which several,
from ATLAS, CMS and LHCb are already available. Specifically, the
inclusion of more processes is envisaged,
which are not currently part of the
dataset, which have a large potential impact on PDFs, and for which higher order corrections have become
available . These include prompt photon
production~\cite{Aaboud:2017cbm}, for which 
NNLO corrections are now available~\cite{Campbell:2016lzl}  and whose
impact on PDF is 
well-known~\cite{d'Enterria:2012yj};  single-top production
(also known at NNLO~\cite{Brucherseifer:2014ama});  and possibly forward  
$D$ meson production or more in general processes with final-state $D$
mesons, such as $W+D$ production, recently measured by
ATLAS~\cite{Aad:2014xca}, a process whose impact on PDFs has been
repeatedly emphasized~\cite{Zenaiev:2015rfa,Gauld:2015yia,Gauld:2016kpd}.

Such a further increase in dataset is likely to require substantial
methodological improvements in PDF determination. Also, it is likely
to result in a further reduction of PDF uncertainties, thereby
requiring better control of theoretical uncertainties. We
specifically expect significant progress in two different
directions. On the one hand, electroweak corrections, which are now
not included, and whose impact is kept under control through kinematic
cuts, will have to be included in a more systematic way. On the other hand,
theoretical uncertainties due to missing higher-order corrections 
will also have to be estimated. Indeed, the preliminary estimates presented in
Sect.~\ref{sec:thunc} suggest that missing higher-order uncertainties
on PDFs, currently not included in the PDF uncertainty, 
are likely to soon become non-negligible, and possibly dominant in
some kinematic regions. All of these improvements will be part of a future
major PDF release. 

\bigskip
\bigskip
\begin{center}
\rule{5cm}{.1pt}
\end{center}
\bigskip
\bigskip

